\section{The exact formula emitted by the new source}
\label{sec:native}

The native numerical constructors are
\[
 \cI(x,y)=4P(x,y),\quad \cR(x,y)=4P(x,y)+1,\quad
 \cQ(x,y)=4P(x,y)+2,\quad \cE(x,y)=4P(x,y)+3.
\]
They encode implication, membership, universal quantification, and
primitive equality. The first argument of $\cQ$ is a variable label.
Set $\cF=1=\cR(0,0)$ and $\cN(t)=\cI(t,1)$. The usual negative
encodings are $A(x,y)=\cN(\cI(x,\cN(y)))$ for conjunction and
$X_i(x)=\cN(\cQ(i,\cN(x)))$ for existence.

For every enabled triple $(a,b,c)$, define the literal expressions
\begin{align}
 r&=\cR(b,c),\nonumber\\
 s&=\cQ(c,\cI(\cR(c,a),\cR(c,b))),\nonumber\\
 w&=\cQ\bigl(b,\cI(c,\cN(\cQ(b,\cI(c,\cN(r)))))\bigr),\label{eq:literal-w}\\
 u&=\cQ(a,\cI(s,w)),\nonumber\\
 h&=\cQ\bigl(b,\cI(r,\cN(\cQ(a,\cI(c,\cN(u)))))\bigr),\label{eq:literal-h}\\
 v&=\cI(\cR(a,b),\cQ(c,\cR(a,b))),\nonumber\\
 \ell&=\cQ(c,\cI(h,\cN(w))).\nonumber
\end{align}
The new gated bundle is exactly
\begin{equation}
 B_{\mathrm{sG}}(a,b,c)=\cN\bigl(\cI(v,\ell)\bigr),
 \qquad c\ne a,b.
 \label{eq:finalbundle}
\end{equation}
For a disabled triple the emitter returns $0$. The new equality row is
exactly
\begin{equation}
 \cI\bigl(\cE(c,a),\cI(\cE(c,b),
                \cQ(\cR(a,b),\cE(a,b)))\bigr).
 \label{eq:finalgcode}
\end{equation}
These are equations between arithmetic expression trees. In particular,
the opaque matrix input is the literal number $c$ and the private
container binder is that same literal number $c$.

For distinct active roles $a\ne b$, after correct decoding and classical logical conversions,
\eqref{eq:finalbundle} expresses
\[
 V_R(a,b,c)\ \land\ \exists C\,(H(C)\land W(a,C)).
\]
The ordering matters to the compiled source even though conjunction
order is immaterial to classical provability. The two occurrences of
$w$, including the one inside $h$, are literally the same expression.
The formerly emitted helper
$e=\cI(\cE(a,b),\cQ(\cR(a,b),\cE(a,b)))$
has disappeared from the old seed bundle $\cN(\cI(v,\cI(e,\ell)))$.

\subsection{Exceptional codes and coincident active labels}

The pinned total decoder $D$ does not treat every numerical implication
uniformly. It sends $0$ to a closed tautology, sends $1$ to $R(0,0)$,
and interprets a native implication with consequent $1$ as genuine named
negation. Other implications use the ordinary decoding identity subject
to the exceptions at consequent $1$ and the pair $(0,0)$.

The new source checker verifies the relevant side conditions uniformly
on its exact expression trees. Intentional negations have consequent
$1$. Every other consequent is a quantifier, an equality, a positive
compound, or a membership atom known different from $1$ by the gate.
The only potentially exceptional membership atoms in that position
involve $b,c$ or $c,b$; the gate excludes their simultaneous zero case.
At least one argument of every genuine implication is positive, so its
code is not $0$. These finite structural arguments cover all natural
parameter values and are not an enumeration of small codes.

The inherited freshness statement is
\[
                    c\notin\FV(D(c))\qquad(c\in\mathbb N).
\]
For positive codes, all occurring names are smaller than the code. The
special decoded formula at $c=0$ is closed, although it contains a bound
occurrence of name $0$. At $c=1$ the only name is $0<1$. Consequently
the actual private label is fresh even in those exceptional cases.

When $a\ne b$, the inherited named-to-two-role abstraction identifies
the decoded formula with the ordinary seed schema. When $a=b$, that
abstraction cannot be used. Write $\psi(x)$ for the decoded matrix with
its sole active name given value $x$. Renaming locally bound variables
for readability gives
\begin{align*}
 W(C)&=(\exists x\,\psi(x))\to
                   \exists x\,(\psi(x)\land R(x,C)),\\
 H(C)&=\forall u\,(R(u,C)\to
                  ((\exists x\,\psi(x))\land W(C))).
\end{align*}
The subset test has equal endpoints and is a tautology; rebinding the
single active name shadows the outer member in the matrix positions.
If no $\psi$ witness exists, an empty $C$ works. If a witness exists,
its singleton works. These objects exist in the inherited upper
interpretation. The split is on existence of a witness, not on the
matrix's value at the old environment. This is the seed source's
existing diagonal argument, unchanged by the independent $G$ row.

\subsection{The lower comparison keeps the matrix-code coupling}

The restricted equality helper supports the inherited interpretation
of genuine equality by membership indiscernibility. For a safe
auxiliary $z=\cE(a,b)>a,b$, put
\[
 E_R(a,b)=
 [\forall z\,(R(z,a)\leftrightarrow R(z,b))]\land
 [\forall z\,(R(a,z)\leftrightarrow R(b,z))].
\]
The translation $\theta$ replaces primitive equality by this formula,
sends genuine bottom to $F$, and preserves membership, implication and
universal binders. The inherited logical bootstrap proves the
translations of the full equality calculus using safe-auxiliary
transport, membership vacuity, and $H_\kappa$. It handles coincidences
of names and does not assume that the chosen numerical auxiliary
commutes literally with arbitrary substitution. Those are substantive
inherited proofs, retained in the predecessor material, not consequences
of the new machine's runtime checks.

For an arbitrary ordinary named seed matrix $f$, first form
\[
 \psi=\theta(f),\qquad p^+=\neg\neg\psi,
 \qquad c=\operatorname{enc}(p^+).
\]
Choose input name $1$ and output name $3$ in the ordinary comparison
syntax. Padding gives $c>3$, and the code-label induction makes $c$
fresh after all equality-translation auxiliaries have been introduced.
Now choose the ordinary seed instance whose matrix is $\neg\neg f$
and whose private name is this $c$. The outer seed template uses only
membership and logical constructors, so its translation has matrix
$\neg\neg\theta(f)$ and private label $c$. Its matrix code is exactly
$c$. Thus its native encoding is the actual emitted instance at
$(1,3,c)$, up to the explicitly derived conjunction and double-negation
conversions. The gate is enabled.

Choosing the code \emph{after} translating equality prevents a
matrix/private-label mismatch. Theorem~\ref{thm:deductive} first
recovers precisely the old seed calculus, so this established lower
comparison can be used without redefining its equality translation.

\section{Construction of the 208-subtree source}
\label{sec:source}

The deterministic source builder starts from the exact seed223 source
and the pinned parent225 macro layout. It makes two semantic changes:
the equality row becomes~\eqref{eq:finalgcode}, and the helper is removed
from the seed bundle as in~\eqref{eq:finalbundle}. It then applies one
closed reassociation. No optimization search is required to reproduce
the released object.

In the old complete instruction word, the equality row occupies the
half-open interval $[486,786)$ and has 300 primitive occurrences. Its
replacement has 409. The seed bundle occupies $[2541,4781)$ and has
2,240 occurrences; its replacement has 1,913. Therefore
\[
                  4962+(409-300)+(1913-2240)=4744.
\]
The extra equality quantifier costs 109 occurrences. Removing the
separate helper and a conjunction constructor saves 327. Every other
primitive occurrence remains in its original order. The new bundle's
postfix word has 65 macro invocations, down from 77.

The checker symbolically executes both the old and new macro words and
compares their expression trees with independently written formulas.
It verifies nine macro contracts against the pinned
\code{Dag230.v}/\code{ProgLib230.v} programs by finite bisimulation for
arbitrary natural register vectors. These include the three parameter
pushes, the constructors, and the push-one constituent. The proof
contracts preserve the prior stack under the emitted formula and require
the inherited clean scratch-register interface.

The initial combined source has a consecutive run of seven zero bits in
one source address. The existing reset routine supports the selected
six-zero invariant, and the predecessor investigations showed that this
constraint matters to actual runtime behavior. We apply the closed right
rotation
\[
                  ((A,B),C)\longmapsto (A,(B,C))
\]
at source address \code{1110100100111000}. Each factor is a closed
source program, meaning its break/continue effects do not escape the
factor. Both independent guarded control-flow resolvers agree before
and after the rotation. The entire primitive word is unchanged. The
rotation is the same local repair used in the merged-equality predecessor;
here it remains compatible with the seed composition.

The result has 208 distinct subtrees: 185 interior nodes and 23 primitive
leaf kinds. Its maximum address length is 36 and maximum zero run is 6.
The independent reference-program comparison checks 6,167 guarded
control-flow cases. An additional finite all-register bisimulation has
1,176 paired boundaries; the inherited seed reconstruction has 1,240.
The exported source receipt records the exact input hashes and formulas.

\section{Compilation and two exact state reductions}
\label{sec:folds}

The machine builder is independent of the formula-emission checker.
It reads the pinned K224 source and transition table and the independently
certified final source. It matches unchanged subtrees, deterministically
assigns compatible old decoder labels to changed nodes, and protects
every state used inside an inherited decoder helper path. It checks
every compound prefix occurrence and the exact state-to-subtree relation.
There are 4,743 such occurrences, one for each internal occurrence in
the full 4,744-leaf binary tree.

The first output is a \emph{new} 224-state seed-plus-$G$ parent, not
the old K224 machine or its old theory. Its K-labelled table has digest
\begin{center}\small\ttfamily
e287d8685bcd75a93d8f01e24b9217381f487a7a970183c8fc2dbec439cc01f8.
\end{center}
Two helper states initially retain no independent source role. The next
two identifications let those helpers implement the two remaining
source roles. All labels in this section are historical K labels; start
is $0$ and halt is $254$.

\subsection{Fold 1: state 185 into state 118}

The only incoming transition to state 159 in the parent is
\[
       \delta(152,0)=(0,L,159).
\]
The only incoming transition to state 118 is
\[
       \delta(159,1)=(1,R,118).
\]
If $\tau$ denotes the tape and $h$ the head position, the assertions
\begin{equation}
 q=159\ \Longrightarrow\ \tau(h+1)=0,
 \qquad q=118\ \Longrightarrow\ \tau(h)=0
 \label{eq:foldinv}
\end{equation}
form an inductive invariant from start state $0$ on every binary tape.
The first incoming transition writes the known zero immediately to the
right of its new head. The second moves onto that known zero. Neither
state is the start state. The independent checker verifies all literal
rows and their target obligations; it does not infer the invariant from
a sampled trace.

The relevant old rows are
\begin{center}
\begin{tabular}{ccl}
\toprule
State & Read & Action\\
\midrule
118 & 0 & $(0,R,155)$\\
118 & 1 & $(1,R,173)$, excluded by the old invariant\\
185 & 0 & $(0,R,155)$\\
185 & 1 & $(1,R,153)$\\
\bottomrule
\end{tabular}
\end{center}
Redirect targets 185 to 118, delete state 185, and replace the one row
of 118 by $(1,R,153)$. Its zero row remains unchanged. Each of the
three possible old branches now has exactly the same write and move,
and the appropriately renamed target, in one new step.

The assertion that 118 scans zero belongs to the \emph{old} parent
machine. It must not be asserted for the folded machine: the new one
row legitimately handles executions formerly in state 185. The proof
relates old and new configurations at simulation boundaries and uses
the invariant only on the old side.

\subsection{Fold 2: state 174 into state 159}

After the first fold, states 174 and 159 have identical rows:
\[
        \delta(q,0)=(0,R,166),\qquad
        \delta(q,1)=(1,R,118),\qquad q\in\{174,159\}.
\]
Redirect 174 to 159 and delete it. This is a one-step quotient with no
tape invariant required.

\begin{proposition}[All-tape correctness of the two folds]
For every binary tape $\tau:\mathbb Z\to\{0,1\}$ and every head
position $h$, the 224-state parent halts from $(0,\tau,h)$ if and only
if the final 222-state machine in historical labels does.
\end{proposition}
\begin{proof}
For the first fold use~\eqref{eq:foldinv}, map state 185 to 118, and fix
all other states. Each allowed old transition is simulated by one new
transition with the same entire tape and head position. All ordinary
rows have only their next state renamed. The second fold has the same
property unconditionally. Compose the simulations. A finite old halt
maps to a finite new halt. If the old run is infinite, its one-step
images give an infinite deterministic new run; it cannot also halt.
\end{proof}

The independent fold checker accepts the supplied certificates and
rejects three negative controls: an incorrect scanned-bit assertion,
omission of the supporting invariant at 159, and the wrong retained
one-row target. These checks establish the specific fold theorem;
they are not used as evidence for the separate logical theory.

\subsection{Canonical labels and the literal count}

Breadth-first traversal from the start, visiting the read-zero edge
before the read-one edge, produces canonical states $0,\ldots,221$.
The unique missing target becomes halt $222$. The maps to historical
K and legacy checker labels are included and checked row by row.
All declared working states are reachable in the directed transition
graph. This is a graph statement, not a claim that every state must
occur in the actual blank-tape execution.

The final decoder has 173 states for 185 interior source nodes, retaining
twelve contextual aliases. The remaining engine and initialization
states number 49. Thus the actual table realizes
\[
                49+185-12=222
\]
working states. The final count is verified directly from the complete
444-row table and its reachable-state set.
